\originalchapter{多极展开与无迹对称张量}
\originalsection{局域源的远场}
设电荷分布支撑在半径 $a$ 的球内, 观察点 $r>a$. 用 $\rr'=x'_i\hat{\vct e}_i$, 核函数的展开为
\[
\frac1{|\rr-\rr'|}=\frac1r+\frac{\nn\cdot\rr'}{r^2}
+\frac{3(\nn\cdot\rr')^2-r'^2}{2r^3}+O(a^3/r^4).
\]
这是把 $(1-2\nn\cdot\rr'/r+r'^2/r^2)^{-1/2}$ 作泰勒展开. 在 $r$ 远大于源尺度时, 高阶项逐渐受抑制. 定义总电荷 $Q=\int\rho\dd^3x'$, 电偶极矩 $\pp=\int\rho\rr'\dd^3x'$, 以及
\[
Q_{ij}=\int\rho(\rr')(3x'_ix'_j-r'^2\delta_{ij})\dd^3x'.
\]
$Q_{ij}=Q_{ji}$ 且 $Q_{ii}=0$, 所以五个独立分量而非九个. 电势远场为
\[
V=\frac1{4\pi\eps}\left[\frac Qr+\frac{\pp\cdot\nn}{r^2}
+\frac{Q_{ij}n_in_j}{2r^3}+\cdots\right].
\]
本稿用带系数3的四极矩定义, 后文所有四极系数依此核对; 文献中另一个归一化不能直接代入.

若平移原点 $\rr''=\rr'-\vct a$, 则 $\pp''=\pp-Q\vct a$. 因而净电荷为零时偶极矩与原点无关; 净电荷不为零时只说“偶极矩是多少”却不指定原点是不完整的.

\originalsection{张量形式与球谐形式}
一个齐次多项式 $H_\ell(\rr)=C_{i_1\cdots i_\ell}x_{i_1}\cdots x_{i_\ell}$ 可取系数张量完全对称, 因为反对称部分与坐标乘积缩并为零. 直接求导给
\[
\nabla^2H_\ell=\ell(\ell-1)C_{jj i_3\cdots i_\ell}x_{i_3}\cdots x_{i_\ell}.
\]
所以无迹正好对应调和齐次多项式. 写 $H_\ell=r^\ell F_\ell(\nn)$ 并代入球坐标拉普拉斯算子, 得 $\Delta_\Omega F_\ell=-\ell(\ell+1)F_\ell$. 因此 $F_\ell$ 是同一阶球谐函数的线性组合. 三维 $\ell$ 阶齐次多项式空间维数为 $\binom{\ell+2}2$, 去掉形如 $r^2P_{\ell-2}$ 的迹部分, 其中 $P_{\ell-2}$ 是任意 $\ell-2$ 阶齐次多项式, 其维数 $\binom\ell2$, 调和部分的维数为 $2\ell+1$ (低阶直接取同一结果). 这与 $m=-\ell,\ldots,\ell$ 的球谐个数一致. 两种写法表达同一空间, 张量写法擅长转动与收缩, 球谐写法擅长边界展开和正交积分.

以二阶为例, 从 $z^2$ 减去其迹得到 $z^2-r^2/3$, 限制在单位球面上并乘 $3/2$ 就是 $P_2(\cos\theta)=(3\cos^2\theta-1)/2$. 三阶单位向量积的无迹部分为
\[
n_{\langle i}n_jn_{k\rangle}=n_in_jn_k-\frac15(n_i\delta_{jk}+n_j\delta_{ik}+n_k\delta_{ij}).
\]
求迹后第一项给 $n_k$, 括号给 $5n_k$, 故系数必须是 $1/5$.

\originalsection{理想偶极的接触项}
在 $r\ne0$ 处微分 $V_\mathrm{dip}=\pp\cdot\nn/(4\pi\eps r^2)$ 得熟悉的 $r^{-3}$ 场. 但该场在原点不是普通可积函数. 精确分布恒等式为
\[
\partial_i\frac{x_j}{r^3}=\PV\frac{\delta_{ij}-3n_in_j}{r^3}+\frac{4\pi}3\delta_{ij}\delta^3(\rr),
\]
其中 $\PV$ 指挖去同心小球后取极限. 对测试函数分部积分, 内球面的项为 $\int n_in_j\varphi(a\nn)\dd\Omega$, 极限等于 $(4\pi/3)\delta_{ij}\varphi(0)$. 其余积分的常数部分在角积分中消失, 测试函数的一阶余项使径向积分收敛. 因而
\[
\ee_\mathrm{dip}=\PV\frac{3(\pp\cdot\nn)\nn-\pp}{4\pi\eps r^3}
-\frac{\pp}{3\eps}\delta^3(\rr).
\]
接触项虽不改变源外场, 却改变包含原点的体积分. 删除它会破坏泊松方程及体平均恒等式. 理想偶极的电荷密度为 $\rho=-\pp\cdot\nabla\delta^3(\rr)$, 这是近邻正负电荷在固定偶极矩极限下的分布.

本章对应 8.07 第8、9讲以及作业5、7的无迹张量和分布题. 张量与球谐的关系在这里从拉普拉斯算子建立, 不将名称对应当作证明.
